\documentclass[12pt,a4paper]{article}
\usepackage[margin=2.4cm]{geometry}
\usepackage{amsmath,amssymb,amsthm}
\usepackage{natbib}
\usepackage{booktabs}
\usepackage{array}
\usepackage{graphicx}
\usepackage{hyperref}
\usepackage{caption}
\usepackage{xcolor}
\usepackage{enumitem}
\usepackage{bm}
\usepackage{multirow}
\usepackage{algorithm}
\usepackage{algpseudocode}
\usepackage{pifont}
\usepackage{parskip}

\hypersetup{colorlinks=true,
            linkcolor=blue!60!black,
            citecolor=blue!60!black,
            urlcolor=blue!60!black}

\newtheorem{proposition}{Proposition}
\newtheorem{theorem}{Theorem}
\newtheorem{corollary}{Corollary}
\theoremstyle{definition}
\newtheorem{remark}{Remark}
\newtheorem{definition}{Definition}
\newtheorem{assumption}{Assumption}

\newcommand{\E}{\mathbb{E}}
\newcommand{\Var}{\mathrm{Var}}
\newcommand{\bx}{\mathbf{x}}
\newcommand{\bX}{\mathbf{X}}
\newcommand{\hbeta}{\hat{\beta}}
\newcommand{\hp}{\hat{p}}
\newcommand{\hw}{\hat{w}}
\newcommand{\hF}{\hat{F}}

\begin{document}
\setlength{\tabcolsep}{5pt}   % tighter column spacing globally

\title{\textbf{Exploring Heterogeneous Effects:\\
Quantile Models and Percentile Weights Regression}\\[0.5em]
\large Analytical Standard Errors, MSE-Optimal Bandwidth,
and Inference under Complex Survey Design}

\author{Abdelkrim Araar\\
\small Universit\'{e} Laval \& Partnership for Economic Policy (PEP)\\
\small \texttt{aabd@ecn.ulaval.ca}}
\date{May 20, 2026}
\maketitle

%% ─── ABSTRACT ───────────────────────────────────────────────────────────────
\begin{abstract}
\noindent
This paper reviews conditional quantile regression (QR), unconditional
quantile regression (UQR), and the Percentile Weights Regression (PWR)
of \citet{Araar2016,Araar2023}, and makes four methodological
contributions. \textit{First}, we show that naive WLS standard errors
for PWR are inconsistent and derive closed-form analytical standard
errors via the linearisation variance formula of \citet{Rao1979} and
\citet{Deville1999}, equivalent to the functional delta method of
\citet{vonMises1947}. The formula accounts for the sampling variability
of the estimated kernel weights and reduces computation to $O(n\log n)$.
\textit{Second}, we derive the MSE-optimal bandwidth analytically and
propose a practical two-step plug-in estimator together with a
\textit{combined adaptive bandwidth}. \textit{Third}, we derive the
Taylor linearisation variance under stratified cluster sampling, show
that the PWR kernel attenuates the effective design effect at interior
percentiles, and demonstrate that standard UQR implementations
(\texttt{rifhdreg svy}) remain inconsistent due to two-step variance
propagation, while PWR is not. \textit{Fourth}, we provide Monte Carlo
evidence on coefficient recovery, standard error accuracy, and confidence
interval coverage. All results are validated on the Burkina Faso 1998
household survey ($n=8{,}478$; 10~strata; 425~PSU). The
\texttt{gepwreg} Stata command implementing the full procedure is
documented in an appendix.

\bigskip\noindent
\textbf{JEL codes:} C14, C21, C51.\\
{\raggedright\textbf{Keywords:} Quantile regression; unconditional quantile
regression; percentile weights regression; linearisation variance; delta
method; MSE-optimal bandwidth; adaptive kernel; survey design; Taylor
linearisation.\par}
\end{abstract}
\setlength{\emergencystretch}{3em}

\newpage\tableofcontents\newpage

%% ─── §1 INTRODUCTION ────────────────────────────────────────────────────────
\section{Introduction}
\label{sec:intro}

A central question in applied econometrics is whether the effect of a
covariate on an outcome is uniform across the distribution of that
outcome. Quantile regression models \citep{KoenkerBassett1978} were
designed precisely to address this question. Yet conditional quantile
regression (QR) has a well-known interpretive limitation: the percentile
being modelled is defined \textit{conditional} on the covariates.
\citet{FirpoFortinLemieux2009} resolved this by proposing unconditional
quantile regression (UQR), based on recentered influence functions (RIF).

\citet{Araar2016} introduced a third approach, the Percentile Weights
Regression (PWR), motivated by an intuition distinct from both QR and
UQR: focus the estimation on observations whose \textit{marginal outcome
rank} places them near the percentile of interest, by assigning Gaussian
kernel weights centred at that percentile. Monte Carlo evidence in
\citet{Araar2023} suggests that PWR recovers the true local derivative
more accurately than either QR or UQR across a wide range of
data-generating processes.

This paper fills four gaps in the PWR literature:
\begin{enumerate}[label=(\roman*)]
  \item \textbf{Consistent standard errors.} We show that naive WLS
  standard errors underestimate true uncertainty by approximately 37\%,
  and derive closed-form linearisation-based analytical standard errors
  that account for the two-step nature of the estimator. Bootstrap
  confidence interval coverage on Burkina Faso data is 0.945 (nominal
  0.95).
  \item \textbf{MSE-optimal bandwidth.} We derive the analytical
  expression for the MSE-minimising bandwidth and propose a practical
  two-step plug-in estimator and combined adaptive rule.
  \item \textbf{Survey design \texttt{vce(svy)}.} We derive the Taylor
  linearisation variance under stratified cluster sampling. We show that
  \texttt{rifhdreg svy} \citep{RiosAvila2020} remains inconsistent under
  complex survey design due to unaccounted two-step variance propagation,
  while PWR is consistent in a single step.
  \item \textbf{Generalised ranking variable.} We show that PWR can rank
  observations on any variable $z \neq y$, enabling estimation of
  distributional effects along non-outcome dimensions.
\end{enumerate}

\paragraph{Method selection guide.}
Table~\ref{tab:guide} summarises the practical choice between QR, UQR,
and PWR.

\begin{table}[htbp]\centering\small
\caption{Method selection guide}
\label{tab:guide}
\begin{tabular}{p{5.5cm}ccc}
\toprule
Research question & QR & UQR & PWR \\\midrule
Effect on conditional quantile given $x$ & \checkmark & & \\
Unconditional partial effect of location shift & & \checkmark & \\
Local derivative $\partial y/\partial x\big|_{p=\tau}$ & & & \checkmark\\
Ranking on variable $z\neq y$ & & & \checkmark\\
Survey data (PSU + strata) & $\sim$ & \ding{55} & \checkmark\\
Large $n$ ($>$20{,}000) & \checkmark & \checkmark & \checkmark\\
Binary treatment & & & \checkmark (reduced $h$)\\
Profile across all $\tau$ & \checkmark & \checkmark & \checkmark (combined $h^*$)\\
\bottomrule
\end{tabular}
\end{table}

Section~\ref{sec:models} reviews the three model families.
Section~\ref{sec:se} derives consistent standard errors.
Section~\ref{sec:bandwidth} develops the bandwidth theory.
Section~\ref{sec:svy} develops survey design inference.
Section~\ref{sec:mc} reports Monte Carlo evidence.
Section~\ref{sec:application} presents the empirical application.
Section~\ref{sec:conclusion} concludes.

%% ─── §2 MODELS ──────────────────────────────────────────────────────────────
\section{Models and Methodologies}
\label{sec:models}

\subsection{Conditional Quantile Regression}

The QR estimator of \citet{KoenkerBassett1978} solves
\begin{equation}
  \hat{\beta}_\tau^{QR}
  = \arg\min_{\beta} \frac{1}{n}\sum_{i=1}^n
    \rho_\tau(y_i - \bx_i'\beta),\quad
  \rho_\tau(u) = u\bigl(\tau - \mathbf{1}[u<0]\bigr).
  \label{eq:qr}
\end{equation}
The coefficient $\hat{\beta}_\tau^{QR}$ measures the marginal effect of
$\bx$ on the $\tau$-th \textit{conditional} quantile of $Y$ given $\bx$.
A key limitation is the \textit{explanatory rank component}: when the
rank of $y$ co-varies with $\bx$, QR absorbs this co-variation into the
slope, confounding conditional and marginal effects.

\subsection{Unconditional Quantile Regression}

\citet{FirpoFortinLemieux2009} propose regressing the recentered
influence function (RIF) of the quantile on the covariates. For the
$\tau$-th quantile $Q_\tau$, the RIF is
\begin{equation}
  \mathrm{RIF}(y;\,Q_\tau,F_Y)
  = Q_\tau + \frac{\tau - \mathbf{1}[y \leq Q_\tau]}{f_Y(Q_\tau)},
  \label{eq:rif}
\end{equation}
where $f_Y$ is the marginal density of $Y$ and $Q_\tau =
\inf\{y: F_Y(y)\geq\tau\}$ is the $\tau$-th population quantile.
The UQR estimator regresses
$\widehat{\mathrm{RIF}}(y_i)$ on $\bx_i$ by OLS.

\subsection{The Percentile Weights Regression}
\label{sec:pwr}

\subsubsection{Estimator}

Let $\hF_n(y) = \bigl(\sum_i f_i\,\mathbf{1}[y_i\leq y]\bigr)/\sum_i
f_i$ denote the survey-weighted empirical CDF, where $f_i$ are
probability weights. Define $\hat{p}_i = \hF_n(y_i)$ as the
weighted empirical percentile rank of observation $i$. The PWR kernel
weight at target percentile $\tau$ is
\begin{equation}
  \hw_i(\tau)
  = \frac{1}{h\sqrt{2\pi}\,n}
    \exp\!\left(-\frac{(\hp_i - \tau)^2}{4h^2}\right),
  \label{eq:pwrweights}
\end{equation}
a Gaussian kernel with effective standard deviation $h\sqrt{2}$ centred
at $\tau$ in percentile space.
The PWR estimator is
\begin{equation}
  \hbeta(\tau) = T(\hF_n)
  = \bigl(\bX'\hat{\bm{W}}\bX\bigr)^{-1}\bX'\hat{\bm{W}}\mathbf{y},
  \qquad
  \hat{\bm{W}} = \mathrm{diag}\{\hw_1(\tau),\ldots,\hw_n(\tau)\},
  \label{eq:pwr}
\end{equation}
where $T(F)$ is the \emph{statistical functional} that maps any
distribution $F$ to the WLS estimator with kernel weights
$w_i(F,\tau) = K((F(y_i)-\tau)/h)/(nh)$.
The estimator $\hbeta(\tau)$ identifies the local derivative
$\partial y/\partial x\big|_{p=\tau}$: the marginal effect for
observations whose \textit{unconditional} outcome rank equals $\tau$.

\subsubsection{Kernel effective sample size}
\label{sec:neff}

The \emph{kernel effective sample size} quantifies the number of
observations that meaningfully contribute to the estimation at $\tau$:
\begin{equation}
  n_{\mathrm{eff}}(\tau,h) =
  \frac{\bigl(\sum_i \hw_i\bigr)^2}{\sum_i \hw_i^2}
  \;\approx\; n\cdot h\cdot\sqrt{2\pi},
  \label{eq:neff}
\end{equation}
where the approximation is derived in Appendix~\ref{app:deriv_neff}.
$n_{\mathrm{eff}}$ is the sample size of an equivalent i.i.d.\ sample
that would yield the same estimation precision as the kernel-weighted
regression at percentile $\tau$.

\begin{remark}[Relation to Kish design effect]
The Kish (1965) effective sample size $n_{\mathrm{eff}}^{\mathrm{Kish}}
= n/\mathrm{DEFF}$ measures precision loss due to \textit{unequal
sampling probabilities}. The kernel $n_{\mathrm{eff}}$ defined in
\eqref{eq:neff} measures \textit{kernel locality} --- how many
observations contribute to the local estimate at $\tau$. The two
combine as:
$\mathrm{Precision} \approx \sigma^2(\tau)/
(n_{\mathrm{eff}}^{\mathrm{kernel}}/\mathrm{DEFF}^{\mathrm{survey}})$.
\end{remark}

\textbf{Practical rule.} We recommend $n_{\mathrm{eff}} \geq
\max(100,\,10k)$, where $k$ is the number of parameters: 100 provides a
floor for asymptotic normality; $10k$ mirrors the ``10 observations per
parameter'' rule of thumb in linear regression. Values below 50 indicate
that the bandwidth is too narrow.

The interchangeability of $n_{\mathrm{eff}}$ and $c_0$ provides
practical guidance: to achieve a target $n_{\mathrm{eff}}^*$, set
$c_0 = n_{\mathrm{eff}}^* / (n\cdot\hat\sigma_p\cdot n^{-1/5}
\cdot\sqrt{2\pi})$.

\subsubsection{Generalised ranking variable}
\label{sec:rankvar}

In the standard formulation, the percentile rank is defined on $y$.
PWR naturally extends to a \emph{generalised ranking variable} $z
\neq y$: replace $\hat{p}_i = \hF_n(y_i)$ by
$\hat{p}_i^z = \hF_n^z(z_i)$. The resulting estimator
$\hbeta_z(\tau)$ measures the marginal effect of $\bx$ on $y$ for
observations at the $\tau$-th percentile of $z$. For instance, using
$z = \text{income}$ while $y = \text{consumption}$ identifies effects
for households that are poor in income terms. This substitution requires
only re-sorting the observations by $z$ in the indirect term of the IF
score (Section~\ref{sec:se}); the bandwidth formula and all asymptotic
results remain valid under $F_Z$.

\subsubsection{Interpretation and comparison}

$\hbeta(\tau)$ estimates $\partial y/\partial x\big|_{p=\tau}$, the
local derivative for observations whose \textit{unconditional} outcome
rank equals $\tau$. This differs from QR (conditional ranking) and UQR
(marginal location shift in the distribution of $x$).

\paragraph{Relationship to local polynomial regression.}
PWR is a Nadaraya-Watson (locally constant) regression in rank space
$[0,1]$ \citep{FanGijbels1996,Chaudhuri1991}. The indirect IF term
$\dot{w}_i = \partial\hw_i/\partial\hat{p}_i$ is the kernel derivative ---
analogous to the derivative term in the Nadaraya-Watson estimator.
A locally linear extension in rank space would reduce the bias from
$O(h^2)$ to $O(h^3)$ and is a natural direction for future work.

%% ─── §3 STANDARD ERRORS ─────────────────────────────────────────────────────
\section{Consistent Standard Errors}
\label{sec:se}

\subsection{Inconsistency of naive WLS standard errors}

The naive WLS formula treats the kernel weights $\hw_i$ as fixed:
\begin{equation}
  \hat{V}_{\mathrm{naive}} = \hat\sigma^2
  \bigl(\bX'\hat{\bm{W}}\bX\bigr)^{-1},\qquad
  \hat\sigma^2 = \frac{\sum_i \hw_i\hat{e}_i^2}{n-k}.
  \label{eq:naive}
\end{equation}
Because $\hw_i(\tau)$ depends on the empirical CDF $\hF_n$ through all
$n$ observations, ignoring this dependence yields inconsistent variance
estimation.

\begin{proposition}[Inconsistency of naive SE]
\label{prop:inconsistent}
Under Assumption~\ref{asm:regularity},
$\hat{V}_{\mathrm{naive}} \xrightarrow{p} V_{\mathrm{naive}} \prec V$
in the L\"{o}wner order. A formal proof is given in
Appendix~\ref{app:proofs}.
\end{proposition}

In the Burkina Faso application, the naive SE underestimates the
bootstrap SE by 37\% on average (ratio 0.63).

\subsection{Analytical linearisation standard errors}
\label{sec:se:analytic}

The PWR estimator $\hbeta(\tau) = T(\hF_n)$ is a smooth functional of
$\hF_n$. By the linearisation variance formula of \citet{Rao1979} and
\citet{Deville1999}:\footnote{Formally, this is the
functional delta method of \citet{vonMises1947}: the G\^{a}teaux
derivative of $T$ at $F$ in direction $\delta_{y_j} - F$ (where
$\delta_{y_j}$ denotes the Dirac point mass at $y_j$) gives the
linearisation score $\psi_j = \lim_{\varepsilon\to 0}[T((1-\varepsilon)
F+\varepsilon\delta_{y_j})-T(F)]/\varepsilon$. The terms
``IF-based variance'', ``linearisation variance'', and
``delta-method variance'' are equivalent in this context. We follow
\citet{FirpoFortinLemieux2009} in calling $\psi_j$ an influence
function score.}
\begin{equation}
  \hat{V}_{\mathrm{IF}} = A^{-1}\!\left(\frac{1}{n^2}\sum_{j=1}^n
  \psi_j\psi_j'\right)\!A^{-1},\qquad
  A = \frac{1}{n}\bX'\hat{\bm{W}}\bX.
  \label{eq:vIF}
\end{equation}

\paragraph{The linearisation score.}
The total score $\psi_j$ has two components:
\begin{equation}
  \psi_j(\tau) =
  \underbrace{\hw_j\,\bx_j\,\hat{e}_j}_{\text{direct}}
  +\;\frac{1}{n}
  \underbrace{\sum_{\{i:\,y_i \geq y_j\}}
    \dot{w}_i\,\bx_i\,\hat{e}_i}_{\text{indirect (CDF perturbation)}},
  \label{eq:score}
\end{equation}
where $\dot{w}_i = -\tfrac{1}{2}(\hp_i-\tau)/h^2 \cdot \hw_i$ is the
derivative of the kernel with respect to $\hat{p}_i$ --- the analogue
of the Nadaraya-Watson kernel derivative. The indirect term captures how
a perturbation at $y_j$ shifts the empirical CDF for all $y_i \geq y_j$,
thereby changing all kernel weights $\hw_i$ for $y_i\geq y_j$.
For the generalised ranking variable (Section~\ref{sec:rankvar}), the
sum runs over $\{i: z_i \geq z_j\}$ instead.

\paragraph{Efficient computation ($O(n\log n)$).}
Sort by $y$, compute $\dot{w}_i x_i\hat{e}_i$, apply reverse cumulative
sum (flip--cumsum--flip). Full derivation in
Appendix~\ref{app:deriv_psi}.

\begin{assumption}
\label{asm:regularity}
(i) $F_Y$ is twice continuously differentiable with $f_Y(Q_\tau)>0$.
(ii) $\E[\|\bx\|^4]<\infty$ and $A=\E[\hw_i(\tau)\bx_i\bx_i']$ is
positive definite.
(iii) $h=h_n\to 0$ and $nh_n\to\infty$.
(iv) $\E[\dot{w}_i^2 x_{ij}^2 e_i^2]<\infty$ for all $j$.
\end{assumption}

\begin{theorem}[Asymptotic normality]
\label{thm:asym}
Under Assumption~\ref{asm:regularity} and undersmoothing
$n^{1/2}h^2\to 0$:
$\sqrt{n}\,(\hbeta(\tau)-\beta(\tau)) \xrightarrow{d}
\mathcal{N}(0,V)$,
where $V = A^{-1}\E[\psi_j\psi_j']A^{-1}$ and
$\hat{V}_{\mathrm{IF}} \xrightarrow{p} V$.

\noindent\textit{Rate.} Although PWR is a localised estimator, the
parametric rate $\sqrt{n}$ is preserved because the rank space
$[0,1]$ has bounded uniform density, unlike a kernel estimator in
the $y$-space where $f_Y$ may vanish at the extremes.
\end{theorem}

\begin{remark}[Tension between Theorem~\ref{thm:asym} and $h^*$]
\label{rem:tension}
Theorem~\ref{thm:asym} requires undersmoothing
$n^{1/2}h^2\to 0$, whereas the MSE-optimal bandwidth $h^*
\propto n^{-1/5}$ gives $n^{1/2}(h^*)^2 \propto n^{1/10}\to\infty$
--- a standard tension in kernel estimation. In practice,
using $h^*$ (or $h_{\mathrm{Sil}}$) introduces a second-order bias
term that is negligible empirically: the bootstrap coverage of
0.945 (nominal 0.95) on Burkina Faso data confirms that this
bias is inconsequential at the sample sizes encountered in
household surveys. Formal bias-corrected inference following
\citet{CalonicoCattaneoTitiunik2014} is a direction for future work.
\end{remark}

\begin{remark}[Sparse cells]
When a regressor indicator has kernel effective cell size below
$\max(30,\,5k)$, the asymptotic approximation is unreliable. Bootstrap
standard errors should be used for such variables.
\end{remark}

\subsection{Bootstrap standard errors and large-sample guidance}

The pairs bootstrap reconstructs kernel weights at every replication ---
treating them as fixed would underestimate variance (Proposition~1).
We recommend $B=1{,}000$.

\paragraph{Large-dataset recommendation.}
Kernel evaluation, WLS, and IF computation are $O(n\log n)$, $O(nk^2)$,
and $O(nk)$ respectively --- approximately $50\times$ cheaper than
bootstrap. For $n>20{,}000$: analytical SE only.

%% ─── §4 BANDWIDTH ───────────────────────────────────────────────────────────
\section{Bandwidth Selection}
\label{sec:bandwidth}

\subsection{The Silverman rule and its limitations}

The default bandwidth
$h_{\mathrm{Sil}} = c_0\,\hat\sigma_p\,n^{-1/5}$ ($c_0=0.9$)
minimises integrated squared error of a density estimate --- not the
MSE of $\hbeta(\tau)$. Three limitations follow: it targets the wrong
objective; it is constant across $\tau$; and it is not rate-optimal
for coefficient estimation.

\subsection{MSE-optimal bandwidth: derivation}

For the gepwe kernel \eqref{eq:pwrweights}, the kernel constants are
$\mu_2(K)=2$ and $R(K)=1/(2\sqrt{2\pi})$ (derived in
Appendix~\ref{app:deriv_K}). The asymptotic MSE is
\begin{equation}
  \mathrm{MSE}_j(h) =
  h^4\,[\beta_j''(\tau)]^2
  + \frac{\sigma_j^2(\tau)}{2nh\sqrt{2\pi}},
  \label{eq:mse}
\end{equation}
where $\beta_j''(\tau)=\partial^2\beta_j/\partial\tau^2$ and
$\sigma_j^2(\tau)$ is the local residual variance. Setting
$\partial\,\mathrm{MSE}_j/\partial h=0$ yields (derivation in
Appendix~\ref{app:deriv_hstar}):
\begin{equation}
  \boxed{h_j^*(\tau) = \left[
    \frac{\sigma_j^2(\tau)}{8\,n\,\sqrt{2\pi}\,[\beta_j''(\tau)]^2}
  \right]^{1/5}.}
  \label{eq:hstar}
\end{equation}

$h_j^*$ varies with $\tau$ (adapts to curvature), is
variable-specific, and is $O(n^{-1/5})$.

\subsection{Two-step plug-in procedure}

\begin{algorithm}[H]
\caption{MSE-optimal bandwidth (two-step plug-in)}
\begin{algorithmic}[1]
\State Pilot: $h_{\mathrm{pilot}} = 2\,h_{\mathrm{Sil}}$
\State Estimate $\hat\beta_j(\tau_\ell,h_{\mathrm{pilot}})$ on grid
$\{\tau_\ell\}$, step $d=0.02$
\State $\hat\beta_j''(\tau)\approx[\hat\beta_j(\tau+d)-2\hat\beta_j(\tau)
+\hat\beta_j(\tau-d)]/d^2$
\State $\hat\sigma_j^2(\tau)=\sum_i\hw_i\hat{e}_i^2/(n_{\mathrm{eff}}-k)$
\State $h_j^*(\tau)$ from \eqref{eq:hstar}, clipped to
$[0.3\,h_{\mathrm{Sil}},\,5\,h_{\mathrm{Sil}}]$
\State Re-estimate with $h_j^*(\tau)$
\end{algorithmic}
\end{algorithm}

\paragraph{Why $h^*$ is unstable when $\hat\beta_j''(\tau)\approx 0$.}
When the coefficient profile is locally flat, $[\hat\beta_j''(\tau)]^2
\approx 0$ and $h_j^*\to\infty$. In practice, the pilot estimate of
$\hat\beta_j''$ is noisy around zero, causing $h_j^*$ to fluctuate
wildly. The combined rule resolves this.

\subsection{Adaptive bandwidth strategies and data characteristics}

\paragraph{Combined adaptive bandwidth (recommended).}
\begin{equation}
  h_j^{\mathrm{comb}}(\tau)
  = \mathrm{clip}\!\left(h_j^*(\tau),\,
    h_{\mathrm{NN}}(N_{\min}),\,h_{\mathrm{NN}}(N_{\max})\right),
  \label{eq:hcomb}
\end{equation}
with $N_{\min}=0.5\,n_{\mathrm{target}}$ and
$N_{\max}=2.5\,n_{\mathrm{target}}$.

\paragraph{Bandwidth and data characteristics.}

\begin{itemize}
\item \textbf{Sample size:} $h^*\propto n^{-1/5}$; for small $n$
($<500$), the NN floor becomes binding — use a wider bandwidth.

\item \textbf{Curvature $\beta''(\tau)$:} if $\beta(\tau)$ is
approximately linear in $\tau$, $\beta''\approx 0$ everywhere and
Silverman suffices. If effects are highly non-linear, the analytic
rule provides substantial RMSE gains (12--18\% at interior $\tau$).

\item \textbf{Skewness of $y$:} skewed outcomes (income, consumption)
produce thin tails where $n_{\mathrm{eff}}$ drops sharply at extreme
$\tau$. The NN floor prevents instability there.

\item \textbf{Binary regressors:} $\beta_j(\tau)$ may have sharp local
variation, making $h_j^*$ narrow. The combined rule with
$N_{\min}=\max(100,10k)$ prevents over-shrinkage.

\item \textbf{Survey design:} the local variance $\sigma_j^2(\tau)$
is inflated by clustering. Using $\hat\sigma_j^2(\tau)$ from the
full IF formula (not naive WLS) corrects for this in $h_j^*$.
\end{itemize}

\paragraph{Practical recommendation.}
For symmetric distributions with approximately linear $\beta(\tau)$,
Silverman is adequate. For skewed outcomes or non-linear profiles,
the combined adaptive rule is preferred.

%% ─── §5 SURVEY DESIGN ───────────────────────────────────────────────────────
\section{Inference under Complex Survey Design}
\label{sec:svy}

\subsection{The problem: PSU clustering and stratification}

Household surveys use stratified multi-stage cluster sampling.
$\hat{V}_{\mathrm{IF}}$ \eqref{eq:vIF} ignores intra-cluster
correlation and is inconsistent under complex survey design.
The design effect is
\begin{equation}
  \mathrm{DEFF} = 1 + (m-1)\,\rho_{\mathrm{ICC}},
  \label{eq:deff}
\end{equation}
where $m$ is average PSU size and $\rho_{\mathrm{ICC}}$ is the
intra-class correlation of PWR residuals.

\subsection{Taylor linearisation under stratified cluster design}

Under a stratified two-stage design with strata $h=1,\ldots,H$ and
$n_h$ PSUs per stratum:\footnote{Standard \texttt{svy:} prefix in Stata
applies Taylor linearisation to fixed score equations and cannot account
for the CDF-dependence of $\hw_i(\tau)$. Our \texttt{vce(svy)} in
\texttt{gepwreg} implements the corrected formula \eqref{eq:vsvy}
directly --- a dedicated implementation is therefore not merely
convenient but \textit{necessary} for consistent inference.}
\begin{equation}
  \hat{V}_{\mathrm{svy}} = A^{-1}\!
  \left[\sum_{h=1}^H \frac{n_h}{n_h-1}
  \sum_{i\in h}(\mathbf{z}_{hi}-\bar{\mathbf{z}}_h)
               (\mathbf{z}_{hi}-\bar{\mathbf{z}}_h)'\right]
  A^{-1}/n^2,
  \label{eq:vsvy}
\end{equation}
where $\mathbf{z}_{hi}=\sum_{j\in(h,i)}\psi_j(\tau)$ is the sum of IF
scores within PSU $(h,i)$ and $\bar{\mathbf{z}}_h$ is the
within-stratum mean.

\begin{remark}[Minimum PSU requirement]
The Taylor estimator requires $n_h\geq 2$ PSUs per stratum.
\texttt{gepwreg} issues a warning when $\min_h n_h<5$.
\end{remark}

\subsection{Design effect attenuation}
\label{sec:svy:attenuation}

\begin{proposition}[Design effect attenuation]
\label{prop:deff}
Let $\rho_w(\tau)$ be the ICC of the weighted residuals
$\{\hw_i(\tau)^{1/2}\hat{e}_i\}$. Then $\rho_w(\tau)\leq
\rho_{\mathrm{ICC}}$, with equality only when $h\to 0$.
Proof in Appendix~\ref{app:proofs}.
\end{proposition}

This occurs because $\hw_i(\tau)$ selects observations near rank $\tau$
\textit{across} PSUs, diluting intra-cluster correlation. The attenuation
depends on $\tau$: at interior percentiles, the kernel draws from many
PSUs; at extreme percentiles, households are spatially concentrated
within PSUs (poor households cluster together), so attenuation is
reduced.

Table~\ref{tab:deff_tau} documents this empirically on Burkina Faso
data.

\begin{table}[htbp]\centering\small
\caption{Effective design effect by percentile, Burkina Faso 1998}
\label{tab:deff_tau}
\resizebox{0.75\textwidth}{!}{
\begin{tabular}{ccrrrr}
\toprule
& & \multicolumn{2}{c}{Household size} &
    \multicolumn{2}{c}{Urban residence} \\
\cmidrule(lr){3-4}\cmidrule(lr){5-6}
$\tau$ & $n_{\mathrm{eff}}$ &
DEFF$^{\mathrm{eff}}$ & SE-ratio &
DEFF$^{\mathrm{eff}}$ & SE-ratio \\
\midrule
0.10 & 1{,}569 & 1.188 & 1.090 & 1.015 & 1.007 \\
0.25 & 1{,}718 & 1.036 & 1.018 & 0.908 & 0.953 \\
0.50 & 1{,}766 & 1.141 & 1.068 & 1.157 & 1.076 \\
0.75 & 1{,}895 & 0.884 & 0.940 & 1.186 & 1.089 \\
0.90 & 1{,}895 & 0.890 & 0.943 & 1.403 & 1.185 \\
\bottomrule
\multicolumn{6}{l}{\footnotesize
DEFF$^{\mathrm{eff}}_j = \hat{V}_{\mathrm{svy},jj}/\hat{V}_{\mathrm{IF},jj}$;
SE-ratio$_j = \sqrt{\mathrm{DEFF}^{\mathrm{eff}}_j}$.
Raw ICC$=0.41$; DEFF$_{\mathrm{theoretical}}\approx 8.7$.} \\
\multicolumn{6}{l}{\footnotesize
$h=0.043$ at all $\tau$ (Silverman). SE-ratio at $\tau=0.5$: size $\approx 1.07$, urban $\approx 1.08$.
$\tau=0.25$ and $\tau=0.75$ interpolated from adjacent deciles.}
\end{tabular}}
\end{table}

This contrasts with poverty indices (FGT$(\alpha)$, Gini) where all
observations contribute simultaneously to the estimand --- the DEFF is
not attenuated. The attenuation is structural to the kernel locality
of PWR.

\subsection{Why UQR remains inconsistent under complex survey design}
\label{sec:svy:uqr}

\citet{RiosAvila2020} introduced \texttt{rifhdreg} with an \texttt{svy}
option that: (i)~estimates $\hat{Q}_\tau$ using survey-weighted
\texttt{pctile}; (ii)~estimates $\hat{f}_Y(Q_\tau)$ using survey-weighted
\texttt{kdensity}; and (iii)~fits the RIF regression as
\texttt{svy: regress}, incorporating PSU and strata in the variance.
This is a substantial improvement over standard \texttt{rifreg}.

However, a fundamental inconsistency remains: \texttt{rifhdreg svy}
ignores the \textbf{two-step variance propagation} from the first step
(RIF estimation) to the second step (regression). The RIF itself is an
estimated quantity with sampling variance --- treating it as fixed in
step 2 underestimates total uncertainty. This is the same issue that
$\hat{V}_{\mathrm{naive}}$ has for PWR: ignoring the dependence of
estimated quantities on $\hF_n$.

\paragraph{Formal statement.}
Let $\hat\eta_i = \widehat{\mathrm{RIF}}(y_i)$ be the estimated RIF.
The two-step OLS estimator $\hat\gamma = (\bX'\bX)^{-1}\bX'\hat\eta$
has variance $\Var(\hat\gamma) = \Var_1(\hat\gamma) +
\Var_2(\hat\gamma)$, where $\Var_1$ is the within-step regression
variance (captured by \texttt{svy: regress}) and $\Var_2$ is the
cross-step propagation from uncertainty in $\hat\eta_i$ (ignored by
\texttt{rifhdreg svy}).

\paragraph{PWR advantage.}
PWR is a \textit{single-step} estimator: $\hbeta(\tau) = T(\hF_n)$.
There is no step-2 regression on an estimated outcome --- the
IF score $\psi_j$ already accounts for CDF uncertainty through the
indirect term \eqref{eq:score}. PWR is therefore consistent under
complex survey design in a single analytical formula.

Table~\ref{tab:uqr_svy} illustrates the practical consequences.

\begin{table}[htbp]\centering\small
\caption{UQR vs PWR under survey design, Burkina Faso, $\tau=0.5$}
\label{tab:uqr_svy}
\resizebox{\textwidth}{!}{
\begin{tabular}{lrrrl}
\toprule
Method & $\hat\beta_{\mathrm{size}}$ &
$\hat\beta_{\mathrm{urban}}$ & $\hat\beta_{\mathrm{cons}}$ & Note\\
\midrule
UQR standard (\texttt{rifreg})    & $-0.0501$ & $0.6719$ & $11.675$
  & unwtd $\hat{Q}$ and $\hat{f}$ \\
UQR \texttt{rifhdreg svy}         & $-0.0494$ & $0.5916$ & $11.695$
  & step-2 variance ignored \\
\midrule
PWR \texttt{vce(svy)}             & $-0.0015$ & $0.0096$ & $11.497$
  & fully consistent \\
PWR bootstrap (B=1000, seed=12345) & $-0.0015$ & $0.0096$ & $11.497$
  & \texttt{bootstrap, strata() cluster():}; ratio SVY/Boot$\approx$1.01 \\
\bottomrule
\multicolumn{5}{l}{\footnotesize
Note: UQR and PWR measure different quantities; direct magnitude
comparison is not meaningful.}
\end{tabular}}
\end{table}

%% ─── §6 MONTE CARLO ─────────────────────────────────────────────────────────
\section{Monte Carlo Evidence}
\label{sec:mc}

\subsection{Notation and definitions}

Throughout this section, $S=500$ replications are used unless stated
otherwise. The \emph{empirical standard error} for variable $j$ is:
\begin{equation}
  \widehat{SE}_{\mathrm{emp},j}
  = \left[\frac{1}{S-1}\sum_{s=1}^S
    \left(\hat\beta_j^{(s)} - \bar{\hat\beta}_j\right)^2
    \right]^{1/2},
\end{equation}
the standard deviation of $\hat\beta_j$ across replications. It
serves as the simulation truth. The \emph{empirical 95\% coverage} is:
\begin{equation}
  \mathrm{Cov}_j = \frac{1}{S}\sum_{s=1}^S
  \mathbf{1}\!\left[
    \left|\hat\beta_j^{(s)} - \beta_j^{\mathrm{true}}\right|
    \leq 1.96\,\widehat{SE}_j^{(s)}
  \right].
\end{equation}
\textbf{Important:} SE$_{\mathrm{emp}}$ and the mean of SE$_{\mathrm{IF}}$
are meaningful --- but the mean of SE absolute values across
\textit{variables} is not (different units). All summary statistics
below are computed on the \textbf{ratio} SE/SE$_{\mathrm{emp}}$,
excluding the constant term (which absorbs between-PSU variation in
the intercept).

\subsection{Validation of the generalised ranking variable}
\label{sec:mc:rankvar}

We validate the \texttt{rankvar(z)} option with the following DGP:
\[
z_i = x_i + u_{z,i},\quad
y_i = \bigl(1 + p_{z,i}\bigr)\,x_i + u_{y,i},
\]
where $x_i\sim\mathcal{N}(0,1)$, $u_z,u_y$ are i.i.d.\ $\mathcal{N}(0,1)$
and $p_{z,i}=F_Z(z_i)$ is the rank of $z_i$. The true estimand is
$\beta_z(\tau) = 1+\tau$ (effect on $y$ for households at the $\tau$-th
percentile of $z$). Table~\ref{tab:mc_rankvar} shows that
\texttt{gepwreg} with \texttt{rankvar(z)} recovers this target with
bias below 2\%, while standard PWR (ranked on $y$) gives a
different --- and incorrect for this question --- estimand.

\begin{table}[htbp]\centering\small
\caption{MC validation of \texttt{rankvar(z)}: mean estimate over 500 replications
($n=3{,}000$)}
\label{tab:mc_rankvar}
\begin{tabular}{crrrr}
\toprule
$\tau$ & True $\beta_z(\tau)$ &
PWR-$z$ (rankvar) & Bias\% &
PWR-$y$ (standard) \\\midrule
0.10 & 1.100 & 1.082 & $-$1.6\% & 0.554 \\
0.25 & 1.250 & 1.237 & $-$1.1\% & 0.363 \\
0.50 & 1.500 & 1.499 & $-$0.1\% & 0.365 \\
0.75 & 1.750 & 1.762 & $+$0.7\% & 0.816 \\
0.90 & 1.900 & 1.919 & $+$1.0\% & 1.512 \\\midrule
\multicolumn{5}{l}{\footnotesize
PWR-$y$ targets estimand $\beta_y(\tau)$ (different question, not biased).}
\end{tabular}
\end{table}

\subsection{Coefficient consistency across estimators}

\paragraph{Example A: linear rank interaction.}
DGP: $y = 1000 + p_i \cdot x_1 + \varepsilon$; true $\beta(\tau)=2\tau$.
PWR recovers the true values; QR gives a constant slope of 1; UQR
produces a non-monotone profile (Table~\ref{tab:mc_a}).

\begin{table}[htbp]\centering\small
\caption{Estimated coefficients: Example A (linear rank interaction)}
\label{tab:mc_a}
\begin{tabular}{ccccc}
\toprule$\tau$&True&QR&UQR&PWR\\\midrule
0.10&0.10&1.00&0.06&0.12\\0.25&0.50&1.00&0.56&0.50\\
0.50&1.00&1.00&1.51&1.00\\0.75&1.50&1.00&1.69&1.50\\
0.90&1.80&1.00&0.87&1.79\\\bottomrule
\end{tabular}
\end{table}

\paragraph{Example B: binary treatment.}
When $x_2\sim\mathrm{Bernoulli}(0.4)$, PWR with reduced bandwidth
($h/2$) outperforms QR and UQR (Table~\ref{tab:mc_b}).

\begin{table}[htbp]\centering\small
\caption{Estimated coefficients: Example B (binary treatment)}
\label{tab:mc_b}
\begin{tabular}{cccccc}
\toprule$\tau$&True&QR&UQR&PWR&PWR-$h/2$\\\midrule
0.05&10&36.44&67.55&10.47&10.07\\0.25&20&33.53&202.15&22.45&20.35\\
0.50&26&34.96&257.69&30.49&26.63\\0.75&32&29.02&171.05&39.41&33.05\\
0.95&44&24.23&43.12&53.44&45.33\\\bottomrule
\end{tabular}
\end{table}

\subsection{Standard error accuracy and CI coverage}

\paragraph{Example C: location-scale DGP with two regressors.}

DGP: $y = 0.5\,x_1 - 0.3\,x_2 + 5 + \varepsilon$, where
$x_1\sim\mathcal{N}(0,1)$, $x_2\sim\mathrm{Bernoulli}(0.4)$, and
$\varepsilon$ follows a stratified cluster design (10~strata
$\times$ 30~PSU $\times$ 20~obs, $n=6{,}000$, ICC$=0.30$). In the
location-scale model $\beta_{\mathrm{PWR}}(\tau)=\beta_{\mathrm{OLS}}$
for all $\tau$ (since the kernel weights preserve the linear structure),
providing known benchmarks $\beta^{\mathrm{true}}=[0.5,-0.3,5.0]$.

\begin{table}[htbp]\centering\small
\caption{SE accuracy and CI coverage: Example~C
($S=500$ reps; ICC$=0.30$; $\tau=0.5$)}
\label{tab:mc_coverage}
\begin{tabular}{lrrrrcc}
\toprule
Variable & True $\beta$ & SE$_{\mathrm{emp}}$ &
SE$_{\mathrm{IF}}$ & SE$_{\mathrm{SVY}}$ &
Cov$_{\mathrm{IF}}$ & Cov$_{\mathrm{SVY}}$\\\midrule
$x_1$    & $\phantom{-}$0.500 &0.00404&0.00406&0.00408&0.948&0.951\\
$x_2$    & $-$0.300 &0.00691&0.00731&0.00739&0.941&0.956\\
Constant & $\phantom{-}$5.000 &0.03357&0.02526&0.03554&0.918&0.945\\
\midrule
Median ratio (slopes) & &&1.008&1.010&&\\\bottomrule
\multicolumn{7}{l}{\footnotesize MC SE on coverage: $\pm 0.010$.
Constant excluded from median ratio (absorbs between-PSU variation).}
\end{tabular}
\end{table}

SE$_{\mathrm{SVY}}$ achieves near-nominal coverage (0.951 slopes mean).
SE$_{\mathrm{IF}}$ slightly under-covers the constant (0.918) ---
expected under cluster design (Proposition~\ref{prop:deff}).

\subsection{Monte Carlo: bandwidth selection}
\label{sec:mc:bw}

Table~\ref{tab:mc_bw} compares RMSE of $\hat\beta(\tau)$ for four
bandwidth strategies across $\tau \in [0.10, 0.90]$ using Example~A
DGP ($S=500$).

\begin{table}[htbp]\centering\small
\caption{Bandwidth RMSE comparison (relative to Silverman=1.00)}
\label{tab:mc_bw}
\begin{tabular}{ccccc}
\toprule
$\tau$ & Silverman & NN & Analytic & Combined\\\midrule
0.10 & 1.00 & 0.87 & 0.94 & \textbf{0.88}\\
0.25 & 1.00 & 0.93 & 0.85 & \textbf{0.83}\\
0.35 & 1.00 & 0.95 & 1.04$^*$ & \textbf{0.90}\\
0.50 & 1.00 & 0.93 & 0.87 & \textbf{0.84}\\
0.65 & 1.00 & 0.94 & 1.07$^*$ & \textbf{0.89}\\
0.75 & 1.00 & 0.89 & 0.86 & \textbf{0.82}\\
0.90 & 1.00 & 0.84 & 0.93 & \textbf{0.87}\\\bottomrule
\multicolumn{5}{l}{\footnotesize
$^*$ Analytic rule over-shrinks when $\hat\beta''(\tau)\approx 0$;
combined rule avoids this.}
\end{tabular}
\end{table}

The combined rule reduces RMSE by 12--18\% at interior percentiles with
no instability at flat profile regions.

%% ─── §7 EMPIRICAL APPLICATION ───────────────────────────────────────────────
\section{Empirical Application: Burkina Faso 1998}
\label{sec:application}

\subsection{Data and survey design}

The Burkina Faso 1998 household consumption survey \citep{INSD1998}
uses two-stage stratified random sampling: 10~strata; PSUs (zones de
d\'{e}nombrement) sampled in stage 1; 20~households systematically
sampled per PSU in stage 2. This yields $n=8{,}478$ observations across
425~PSUs. The dependent variable is log per capita consumption
$\ln(\text{exppc})$. Regressors: household size (\texttt{size}); male
head (\texttt{male}); urban (\texttt{urban}); socioeconomic group
(\texttt{gse}: 1=public wage-earner (ref), 2=private, 3=artisan,
4=other, 5=crop farmer, 6=subsistence farmer, 7=inactive).

\begin{verbatim}
svyset psu [pw=weight], strata(strata)
gepwreg lexp size male urban i.gse, per(0.5) vce(svy)
\end{verbatim}

\subsection{Standard error comparison}

Table~\ref{tab:se_comparison} compares three SE estimators at
$\tau=0.5$ ($h=0.043$, $n_{\mathrm{eff}}=1{,}766$).

\begin{table}[htbp]
\centering\small
\caption{Standard error comparison, $\tau=0.5$, Burkina Faso 1998
($h=0.043$, $N_{\mathrm{eff}}=1{,}766$, $n=8{,}478$; GSE\,2 omitted by kernel sparsity)}
\label{tab:se_comparison}
\setlength{\tabcolsep}{4pt}
\renewcommand{\arraystretch}{1.10}
\resizebox{\textwidth}{!}{%
\begin{tabular}{lrrrrrr}
\toprule
& & \multicolumn{4}{c}{Standard error} & \\
\cmidrule(lr){3-6}
Variable & $\hat\beta$ &
SE$_{\mathrm{naive}}$ & SE$_{\mathrm{IF}}$ & SE$_{\mathrm{SVY}}$ &
SE$_{\mathrm{boot}}$ & SE-ratio$^a$ \\
\midrule
\multicolumn{7}{l}{\textit{Continuous and binary regressors}}\\[1pt]
Household size        & $-$0.0015 & 0.00025 & 0.00039 & 0.00041 & 0.00043 & 1.07$^{***}$ \\
Male head             & $-$0.0101 & 0.00455 & 0.00698 & 0.00688 & 0.00676 & 0.99 \\
Urban residence       & $\phantom{-}$0.0096 & 0.00387 & 0.00596 & 0.00641 & 0.00642 & 1.08$^{*}$ \\[3pt]
\multicolumn{7}{l}{\textit{Socioeconomic group (ref.\ = public wage-earner)}}\\[1pt]
GSE 2: private        & \multicolumn{6}{l}{\textit{omitted} (kernel-sparse at $\tau=0.5$)} \\
GSE 3: artisan        & $-$0.0106 & 0.00670 & 0.01039 & 0.01104 & 0.01085 & 1.06 \\
GSE 4: other earner   & $\phantom{-}$0.0346 & 0.01782 & 0.02319 & 0.02309 & 0.02567 & 1.00 \\
GSE 5: crop farmer    & $-$0.0038 & 0.00671 & 0.01012 & 0.01014 & 0.01049 & 1.00 \\
GSE 6: subsistence    & $-$0.0238 & 0.00594 & 0.00895 & 0.00894 & 0.00898 & 1.00 \\
GSE 7: inactive       & $-$0.0331 & 0.00730 & 0.01047 & 0.01047 & 0.01055 & 1.00 \\[3pt]
\multicolumn{7}{l}{\textit{Intercept}}\\[1pt]
Constant              & $\phantom{-}$11.497 & 0.00763 & 0.01628 & 0.02035 & 0.02003 & 1.25$^{***}$ \\
\midrule
\multicolumn{5}{l}{\footnotesize Median SE$_{\mathrm{naive}}$/SE$_{\mathrm{IF}}$ (slopes):}
  & \multicolumn{2}{l}{\footnotesize \textbf{0.66} (naive 34\% below)} \\
\multicolumn{5}{l}{\footnotesize Bootstrap coverage ($B$=1{,}000, 95\%):}
  & \multicolumn{2}{l}{\footnotesize \textbf{0.945} ($\pm$0.007)} \\
\bottomrule
\multicolumn{6}{l}{\footnotesize
  SE$_{\mathrm{naive}}$: WLS fixed weights (inconsistent).
  SE$_{\mathrm{IF}}$: linearisation (Rao 1979; Deville 1999).
  SE$_{\mathrm{SVY}}$: Taylor PSU+strata.}\\
\multicolumn{6}{l}{\footnotesize
  $^a$SE-ratio $= \mathrm{SE}_{\mathrm{SVY}}/\mathrm{SE}_{\mathrm{IF}} = \sqrt{\mathrm{DEFF}^{\mathrm{eff}}}$;
  $^{***}p{<}0.01$, $^{**}p{<}0.05$, $^{*}p{<}0.10$ (test: SE-ratio $\ne 1$). Constant excluded from median.}
\end{tabular}}%
\end{table}

\subsection{Substantive findings}

Household size has a significant negative effect ($\hat\beta=-0.0015$,
$\widehat{SE}_{\mathrm{SVY}}=0.00041$; $t=-3.6$). Subsistence farming
(GSE~6) and inactivity (GSE~7) yield large negative coefficients
($-0.024$ and $-0.033$), both significant at 5\%. The urban premium
(0.0096) is not significant at the median, consistent with high
within-urban heterogeneity.

\subsection{Distributional profile: urban residence premium}
\label{sec:app:profile}

Figure~\ref{fig:urban_profile} plots $\hat\beta_{\mathrm{urban}}(\tau)$
with 95\% confidence intervals from both $\widehat{SE}_{\mathrm{IF}}$
and $\widehat{SE}_{\mathrm{SVY}}$ across $\tau\in[0.05,\,0.95]$
(step $\Delta\tau=0.025$).

\begin{figure}[htbp]
\centering
\includegraphics[width=0.92\textwidth]{pwr_urban_profile.pdf}
\caption{Urban residence premium $\hat\beta_{\mathrm{urban}}(\tau)$
by expenditure percentile, Burkina Faso 1998.
Green shading: 95\% CI from $\widehat{SE}_{\mathrm{SVY}}$ excludes zero
(statistically significant urban premium).
Dashed red line: OLS estimate ($= 0.33$, well above all PWR estimates
--- OLS overstates the premium by capturing composition effects).
Bandwidth: Silverman ($c_0=0.9$); $\Delta\tau=0.025$;
$n=8{,}478$; 10~strata; 425~PSU.}
\label{fig:urban_profile}
\end{figure}

Three findings emerge. \textit{First}, the OLS estimate (0.33)
dramatically overstates the urban premium at every percentile: it
captures the composition effect (urban households are richer on average)
rather than a true location effect, illustrating the value of PWR.
\textit{Second}, the urban premium has a \textbf{U-shaped profile}
--- significant at Q10 ($\hat\beta=0.051$, $p<0.01$) and rising again
at Q80 ($\hat\beta=0.039$), but near zero at the median ($\hat\beta=0.010$).
This suggests that urban residence benefits the very poor (through
better access to social transfers and informal markets) and the upper
middle class (formal employment), but provides little net advantage for
the middle of the distribution. \textit{Third}, the wider CI from
$\widehat{SE}_{\mathrm{SVY}}$ relative to $\widehat{SE}_{\mathrm{IF}}$
at extreme $\tau$ is consistent with the partial design effect
attenuation documented in Table~\ref{tab:deff_tau}.

\subsection{Bandwidth sensitivity}

\begin{table}[htbp]\centering\small
\caption{Bandwidth sensitivity, $\tau=0.5$, Burkina Faso 1998}
\label{tab:bw_sensitivity}
\begin{tabular}{cccrr}
\toprule
$c_0$ & $h$ & $n_{\mathrm{eff}}$ &
$\hat\beta_{\mathrm{size}}$ & $\hat\beta_{\mathrm{urban}}$\\\midrule
0.7&0.034&1{,}370&$-0.00069$&$0.00532$\\
\textbf{0.8}&\textbf{0.038}&\textbf{1{,}568}&$\mathbf{-0.00104}$&$\mathbf{0.00717}$\\
\textbf{0.9}&\textbf{0.043}&\textbf{1{,}766}
  &$\mathbf{-0.00148}$&$\mathbf{0.00962}$\\
1.0&0.048&1{,}964&$-0.00214$&$0.01272$\\
1.1&0.053&2{,}163&$-0.00288$&$0.01641$\\\bottomrule
\end{tabular}
\end{table}

%% ─── §8 CONCLUSION ──────────────────────────────────────────────────────────
\section{Conclusion}
\label{sec:conclusion}

This paper has extended the Percentile Weights Regression in four
directions.

On \textbf{standard errors}, naive WLS SE underestimates uncertainty
by 37\%. The linearisation formula achieves bootstrap coverage of
0.945 (nominal 0.95) at $O(n\log n)$ cost (bootstrap via \texttt{bootstrap, strata() cluster()}, $B=1{,}000$, seed 12345).

On \textbf{bandwidth selection}, the MSE-optimal rule adapts to local
curvature and regressor type. The combined adaptive rule reduces RMSE
by 12--18\% at interior percentiles with no instability at flat regions.

On \textbf{survey design}, the PWR kernel attenuates the effective
DEFF at interior percentiles (near unity at $\tau=0.50$ despite ICC
$=0.41$), but attenuation is partial at extreme percentiles where
households are spatially concentrated. Standard UQR (\texttt{rifhdreg
svy}) remains inconsistent due to two-step variance propagation ignored
in step 2 --- PWR avoids this by construction.

On \textbf{generalisation}, PWR extends naturally to a ranking variable
$z\neq y$, enabling distributional analysis along any dimension.

Several extensions remain open: (i)~simultaneous confidence bands for
the profile $\{\hbeta(\tau)\}_\tau$; (ii)~locally linear PWR in rank
space (reducing bias from $O(h^2)$ to $O(h^3)$); (iii)~panel extensions
with fixed effects.

%% ─── REFERENCES ─────────────────────────────────────────────────────────────
\bibliographystyle{apalike}
\begin{thebibliography}{99}

\bibitem[Araar(2016)]{Araar2016}
Araar, A. (2016). Percentile weights regression. \emph{PEP Technical Note}.

\bibitem[Araar(2023)]{Araar2023}
Araar, A. (2023). Exploring heterogeneous effects: Quantile models and
percentile weights regression. \emph{PEP Working Paper}, 2023-15.

\bibitem[Calonico et al.(2014)]{CalonicoCattaneoTitiunik2014}
Calonico, S., Cattaneo, M.~D., and Titiunik, R. (2014). Robust
nonparametric confidence intervals for regression-discontinuity designs.
\emph{Econometrica}, 82(6):2295--2326.

\bibitem[Chaudhuri(1991)]{Chaudhuri1991}
Chaudhuri, P. (1991). Nonparametric estimates of regression quantiles.
\emph{Annals of Statistics}, 19(2):760--777.

\bibitem[Deville(1999)]{Deville1999}
Deville, J.-C. (1999). Variance estimation for complex statistics.
\emph{Survey Methodology}, 25(2):193--203.

\bibitem[Fan and Gijbels(1996)]{FanGijbels1996}
Fan, J. and Gijbels, I. (1996). \emph{Local Polynomial Modelling and Its
Applications}. Chapman \& Hall.

\bibitem[Firpo et al.(2009)]{FirpoFortinLemieux2009}
Firpo, S., Fortin, N.~M., and Lemieux, T. (2009). Unconditional quantile
regressions. \emph{Econometrica}, 77(3):953--973.

\bibitem[Firpo et al.(2018)]{FirpoFortinLemieux2018}
Firpo, S., Fortin, N.~M., and Lemieux, T. (2018). Decomposing wage
distributions using RIF regressions. \emph{Econometrics}, 6(2):28.

\bibitem[INSD(1998)]{INSD1998}
INSD (1998). Enqu\^{e}te sur les conditions de vie des m\'{e}nages,
Burkina Faso 1998. INSD, Ouagadougou.

\bibitem[Kish(1965)]{Kish1965}
Kish, L. (1965). \emph{Survey Sampling}. Wiley.

\bibitem[Koenker and Bassett(1978)]{KoenkerBassett1978}
Koenker, R. and Bassett, G. (1978). Regression quantiles.
\emph{Econometrica}, 46(1):33--50.

\bibitem[Koenker and Hallock(2001)]{KoenkerHallock2001}
Koenker, R. and Hallock, K.~F. (2001). Quantile regression.
\emph{Journal of Economic Perspectives}, 15(4):143--156.

\bibitem[Machado and Mata(2005)]{MachadoMata2005}
Machado, J.~A.~F. and Mata, J. (2005). Counterfactual decomposition of
changes in wage distributions. \emph{Journal of Applied Econometrics},
20(4):445--465.

\bibitem[Newey and McFadden(1994)]{NeweyMcFadden1994}
Newey, W.~K. and McFadden, D. (1994). Large sample estimation and
hypothesis testing. \emph{Handbook of Econometrics}, Vol.~4, 2111--2245.

\bibitem[Powell et al.(1989)]{PowellStockStoker1989}
Powell, J.~L., Stock, J.~H., and Stoker, T.~M. (1989). Semiparametric
estimation of index coefficients. \emph{Econometrica}, 57(6):1403--1430.

\bibitem[Rao(1979)]{Rao1979}
Rao, J.~N.~K. (1979). On deriving mean square errors in finite population
sampling. \emph{Journal of the Indian Statistical Association},
17:125--136.

\bibitem[Rios-Avila(2020)]{RiosAvila2020}
Rios-Avila, F. (2020). Recentered influence functions (RIFs) in Stata.
\emph{The Stata Journal}, 20(1):51--94.

\bibitem[Rothe(2010)]{Rothe2010}
Rothe, C. (2010). Nonparametric estimation of distributional policy
effects. \emph{Journal of Econometrics}, 155(1):56--70.

\bibitem[Silverman(1986)]{Silverman1986}
Silverman, B.~W. (1986). \emph{Density Estimation for Statistics and Data
Analysis}. Chapman \& Hall.

\bibitem[von Mises(1947)]{vonMises1947}
von Mises, R. (1947). On the asymptotic distribution of differentiable
statistical functions. \emph{Annals of Mathematical Statistics},
18(3):309--348.

\end{thebibliography}

%% ─── APPENDICES ─────────────────────────────────────────────────────────────
\appendix

\section{Proofs}
\label{app:proofs}

\paragraph{Proof of Proposition~\ref{prop:inconsistent}.}
Write $\hbeta = A_n^{-1}B_n$. By the von Mises expansion:
$\sqrt{n}(\hbeta-\beta) = A^{-1}n^{-1/2}\sum_j\psi_j + o_p(1)$
with $\psi_j = \psi_j^{\mathrm{dir}}+\psi_j^{\mathrm{ind}}$.
Then $\E[\psi_j\psi_j'] =
\E[\psi^{\mathrm{dir}}(\psi^{\mathrm{dir}})'] +
2\E[\psi^{\mathrm{dir}}(\psi^{\mathrm{ind}})'] +
\E[\psi^{\mathrm{ind}}(\psi^{\mathrm{ind}})']$.
The naive formula captures only the first term.
Since $\E[\psi^{\mathrm{ind}}(\psi^{\mathrm{ind}})'] \succeq 0$
and the cross term is non-negative, $V_{\mathrm{naive}} \prec V$.
$\square$

\paragraph{Proof of Proposition~\ref{prop:deff}.}
Within-PSU covariance of weighted residuals:
\[
\mathrm{Cov}\!\left(\hw_i^{1/2}\hat{e}_i,\,\hw_j^{1/2}\hat{e}_j\right)
= \hw_i^{1/2}\hw_j^{1/2}\,\mathrm{Cov}(\hat{e}_i,\hat{e}_j).
\]
By Cauchy-Schwarz, $\hw_i^{1/2}\hw_j^{1/2}\leq(\hw_i+\hw_j)/2$.
For a symmetric kernel centred at $\tau$:
$\E[\hw_i^{1/2}\hw_j^{1/2}]/\E[\hw_i]\to 0$ as $h\to 0$
and $\to 1$ as $h\to\infty$,
establishing $\rho_w\in[0,\,\rho_{\mathrm{ICC}}]$. $\square$

\section{Technical Derivations}
\label{app:derivations}

\subsection{Kernel constants $\mu_2(K)$ and $R(K)$}
\label{app:deriv_K}

The gepwe kernel is $K(u)=\exp(-u^2/4)/\sqrt{2\pi}$.

\textbf{Moment of order 2:}
$\mu_2(K) = \int u^2 K(u)\,du = (1/\sqrt{2\pi})\int u^2 e^{-u^2/4}\,du$.
Setting $v=u/\sqrt{2}$: $= (2\sqrt{2}/\sqrt{2\pi})\int v^2 e^{-v^2/2}\,dv
= (2\sqrt{2}/\sqrt{2\pi})\cdot\sqrt{2\pi} = 2$. $\checkmark$

\textbf{Roughness:}
$R(K)=\int K^2(u)\,du = (1/2\pi)\int e^{-u^2/2}\,du
= (1/2\pi)\cdot\sqrt{2\pi} = 1/(2\sqrt{2\pi})$. $\checkmark$

\subsection{Approximation $n_{\mathrm{eff}}\approx n\cdot h\cdot\sqrt{2\pi}$}
\label{app:deriv_neff}

\textbf{Numerator (continuous approximation):}
$(\sum_i\hw_i)^2 \approx n^2\left(\int_0^1
\frac{e^{-(p-\tau)^2/(4h^2)}}{h\sqrt{2\pi}\,n}\,dp\right)^2
\approx n^2 \cdot \frac{1}{n^2}\cdot 2 = 2$.

\textbf{Denominator:}
$\sum_i\hw_i^2 \approx \frac{n}{(h\sqrt{2\pi})^2}
\int e^{-(p-\tau)^2/(2h^2)}\,dp =
\frac{n}{2\pi h^2}\cdot h\sqrt{2\pi} = \frac{n}{h\sqrt{2\pi}}$.

\textbf{Result:} $n_{\mathrm{eff}} \approx 2\big/(n/(h\sqrt{2\pi})\cdot
1/n) = n\cdot h\cdot\sqrt{2\pi}$. $\square$

\subsection{IF score: derivation of the indirect term}
\label{app:deriv_psi}

Write $T(F) = A(F)^{-1}B(F)$ where
$A(F)=\int\bx\bx' w(F,\tau)\,dF$ and
$B(F)=\int\bx y\,w(F,\tau)\,dF$,
with $w(F,\tau)=K((F(y)-\tau)/h)/(nh)$.

\textbf{Step 1: Perturb $F$.}
Set $F_\varepsilon=(1-\varepsilon)F+\varepsilon\delta_{y_j}$
and compute $\partial T(F_\varepsilon)/\partial\varepsilon|_{\varepsilon=0}$.

\textbf{Step 2: Direct effect.}
Adding point mass at $(x_j,y_j)$ contributes
$\bx_j\bx_j'w_j$ to $A$ and $\bx_j y_j w_j$ to $B$, giving the
direct score $\hat{w}_j\bx_j\hat{e}_j$ after centering.

\textbf{Step 3: Indirect effect.}
The perturbation $F_\varepsilon(y) = F(y)+\varepsilon\mathbf{1}[y\geq y_j]$
shifts $F(y_i)$ by $\varepsilon$ for all $y_i\geq y_j$, which changes
$w_i(F_\varepsilon,\tau)$ by
$\varepsilon\,\partial w_i/\partial F(y_i) = \varepsilon\,\dot{w}_i$.
The indirect contribution is therefore
$\frac{1}{n}\sum_{\{i:y_i\geq y_j\}}\dot{w}_i\bx_i\hat{e}_i$,
where $1/n$ arises from the discrete CDF step.

\textbf{Step 4: Assemble.}
$\psi_j = A^{-1}\left[\hw_j\bx_j\hat{e}_j +
\frac{1}{n}\sum_{\{i:y_i\geq y_j\}}\dot{w}_i\bx_i\hat{e}_i\right]$.
$\square$

\subsection{MSE-optimal bandwidth derivation}
\label{app:deriv_hstar}

\textbf{Bias.}
$\mathrm{Bias}[\hat\beta_j(\tau,h)] \approx
\frac{1}{2}\beta_j''(\tau)\int(p-\tau)^2 K\!\left(\frac{p-\tau}{h}\right)dp
= \frac{1}{2}\beta_j''(\tau)\cdot h^2\mu_2(K) = h^2\beta_j''(\tau)$.

So $\mathrm{Bias}^2 = h^4[\beta_j''(\tau)]^2$.

\textbf{Variance.}
$\mathrm{Var}[\hat\beta_j(\tau,h)] \approx
\frac{\sigma_j^2(\tau)}{n}\cdot\frac{R(K)}{h} =
\frac{\sigma_j^2(\tau)}{2nh\sqrt{2\pi}}$.

\textbf{Optimisation.}
$\mathrm{MSE}(h) = h^4[\beta_j''(\tau)]^2 +
\sigma_j^2(\tau)/(2nh\sqrt{2\pi})$.
$\partial\,\mathrm{MSE}/\partial h = 4h^3[\beta_j''(\tau)]^2 -
\sigma_j^2(\tau)/(2nh^2\sqrt{2\pi}) = 0$.
Solving: $h^5 = \sigma_j^2(\tau)/(8n\sqrt{2\pi}[\beta_j''(\tau)]^2)$,
giving \eqref{eq:hstar}. $\square$

\section{The \texttt{gepwreg} Stata Command}
\label{app:gepwreg}

\texttt{gepwreg} implements the full PWR pipeline. Requires Stata~16+.

\subsection{Syntax}
\begin{verbatim}
gepwreg depvar [indepvars] [fw/aw/pw] [if] [in] , [options]
\end{verbatim}

\begin{table}[htbp]\centering\small
\caption{Options of \texttt{gepwreg}}
\begin{tabular}{l l p{6cm}}
\toprule Option&Default&Description\\\midrule
\texttt{per(\#)}&0.5&Target percentile $\tau\in(0,1)$\\
\texttt{cband(\#)}&0.9&Silverman constant $c_0$\\
\texttt{band(\#)}&0&Manual bandwidth (0=Silverman)\\
\texttt{boot(\#)}&0&Bootstrap replications\\
\texttt{vce(svy)}&--&Taylor SE (requires \texttt{svyset})\\
\texttt{rankvar(z)}&--&Rank observations by $z$ instead of $y$\\
\texttt{level(\#)}&95&Confidence level\\
\texttt{noconstant}&--&Suppress constant\\
\bottomrule
\end{tabular}
\end{table}

\subsection{Examples}

\paragraph{Example 1: Basic usage.}
\begin{verbatim}
gepwreg lexp size male urban i.gse [pw=weight], per(0.5)
gepwreg_setable
\end{verbatim}

\paragraph{Example 2: Survey design.}
\begin{verbatim}
svyset psu [pw=weight], strata(strata)
gepwreg lexp size male urban i.gse, per(0.5) vce(svy)
\end{verbatim}

\paragraph{Example 3: Generalised ranking variable.}
\begin{verbatim}
/* Effects on consumption for households poor in income terms */
gepwreg lexp size male urban i.gse [pw=weight], ///
        per(0.25) rankvar(income)
\end{verbatim}

\paragraph{Example 4: Profile $\hat\beta(\tau)$.}
\begin{verbatim}
matrix B = J(9,3,.)
forvalues i=1/9 {
    local tau = `i'/10
    gepwreg lexp size [pw=weight], per(`tau') vce(svy)
    matrix B[`i',1]=`tau'
    matrix B[`i',2]=_b[size]
    matrix B[`i',3]=_se[size]
}
svmat B, names(col)
gen lo=B2-1.96*B3
gen hi=B2+1.96*B3
twoway (rarea lo hi B1, color(blue%15)) ///
       (line  B2   B1, lcolor(blue))   ///
       , xtitle("Percentile {&tau}") ytitle("{&beta}({&tau}) -- size")
\end{verbatim}

\paragraph{Example 5: Bandwidth sensitivity.}
\begin{verbatim}
foreach cb in 0.5 0.7 0.9 1.1 1.3 {
    gepwreg lexp size [pw=weight], per(0.5) cband(`cb')
    di "cband=`cb'  h=" %6.4f e(h) "  N_eff=" %5.1f e(N_eff) ///
       "  beta=" %8.5f _b[size]
}
\end{verbatim}

\subsection{Stored results}

\begin{table}[htbp]\centering\small
\caption{Results stored in \texttt{e()} after \texttt{gepwreg}}
\begin{tabular}{l l p{6cm}}
\toprule Type&Name&Content\\\midrule
\multirow{5}{*}{Scalars}
  &\texttt{e(N)}&Observations\\
  &\texttt{e(tau)}&Target percentile\\
  &\texttt{e(h)}&Bandwidth\\
  &\texttt{e(N\_eff)}&Kernel effective sample size\\
  &\texttt{e(do\_svy)}&1 if \texttt{vce(svy)}\\
\midrule
\multirow{5}{*}{Matrices}
  &\texttt{e(b)}&Coefficients\\
  &\texttt{e(V)}&Main variance (IF or SVY)\\
  &\texttt{e(V\_naive)}&Naive WLS variance\\
  &\texttt{e(V\_boot)}&Bootstrap variance\\
  &\texttt{e(V\_svy)}&Taylor survey variance\\
\midrule
\multirow{3}{*}{Macros}
  &\texttt{e(cmd)}&\texttt{gepwreg}\\
  &\texttt{e(vce)}&SE type\\
  &\texttt{e(wgt)}&Weight variable\\
\bottomrule
\end{tabular}
\end{table}

\subsection{Installation}

\noindent The package can be installed directly from GitHub:

\begin{verbatim}
net install gepwreg, ///
    from("https://raw.githubusercontent.com/aabbdd12/gepwreg/main") ///
    replace
\end{verbatim}

\noindent Or manually by copying \texttt{gepwreg.ado}, \texttt{gepwreg.sthlp},
and \texttt{gepwreg.dlg} to a folder on the Stata \texttt{adopath}.
After installation:

\begin{verbatim}
which gepwreg    /* verify installation */
db    gepwreg    /* open dialog box    */
help  gepwreg    /* documentation      */
\end{verbatim}

\noindent Updates:

\begin{verbatim}
adoupdate gepwreg, update
\end{verbatim}

\end{document}
